% nplayer.tex -- Section 6: N-Player & Social Dilemma Extensions

\section{N-Player Extensions}
\label{sec:nplayer}

While two-player games isolate fundamental strategic tensions, many real-world
alignment challenges arise in multi-agent settings. Kant extends the core
framework to support $N$-player interactions, introducing social dilemma
structures that test collective action, free-riding, and coordination at scale.

% TODO: Add introductory paragraph on why N-player matters for alignment

\subsection{N-Player Games}
\label{sec:nplayer-games}

The following games natively support $N > 2$ players:

\begin{itemize}[nosep]
    \item \textbf{Public Goods Game} --- each of $N$ players contributes from
          a private endowment; the pool is multiplied and split equally.
          Tests free-riding incentives as group size grows.
    \item \textbf{Volunteer's Dilemma} --- at least one player must pay a cost
          to provide a public benefit. Diffusion of responsibility increases
          with $N$.
    \item \textbf{El Farol Bar Problem} --- players independently decide whether
          to attend a venue with limited capacity. Tests minority coordination.
    \item \textbf{Tragedy of the Commons} --- players extract from a shared
          resource; overexploitation depletes it. Tests sustainable cooperation.
\end{itemize}

% TODO: Formal game specifications with N-player payoff functions

\subsection{Scaling from Dyadic to Multi-Agent}
\label{sec:scaling}

Extending from two to $N$ players introduces several challenges:

\begin{itemize}[nosep]
    \item \textbf{Observation space}: agents must reason about $N-1$ opponents
          rather than a single partner, increasing cognitive demands.
    \item \textbf{Action aggregation}: payoffs depend on aggregate opponent
          behavior (e.g., total contributions) rather than individual actions.
    \item \textbf{Strategy complexity}: the opponent strategy space grows
          combinatorially; we address this via role-based strategy assignment.
    \item \textbf{Metric adaptation}: cooperation rate and fairness metrics
          generalize naturally; exploitation resistance requires redefinition
          against adversarial coalitions rather than single defectors.
\end{itemize}

% TODO: Formal metric extensions for N-player settings
% TODO: Experimental results comparing 2-player vs N-player behavior
